% ════════════════════════════════════════════════════════════════════
%  tầng a, phần 3 — c6 · c7 · chọn S*    rồi mở sang tầng b
% ════════════════════════════════════════════════════════════════════

\begin{frame}[fragile]{c6 · chứng minh «không đoạn nào thừa»}
  \framesubtitle{c6}
  \lead{cách hiển nhiên là bỏ từng đoạn một rồi đo lại. cách đó \no{sai}, và
        lý do là $v$ không đơn điệu.}

  \begin{columns}[t]
  \begin{column}{0.47\textwidth}
    \begin{badb}[title={một ca khác — ba đoạn $x$, $y$, $z$}]
      \begin{dia}
tau_q = 0,76
v({x})     = 0,85  >= tau_q   <- x ĐÃ ĐỦ
v({x,y})   = 0,60  <  tau_q      y làm lạc
v({x,z})   = 0,71  <  tau_q
v({y,z})   = 0,40  <  tau_q
v({x,y,z}) = 0,80  >= tau_q

bỏ-một-đoạn từ {x,y,z}:
   ba cặp ĐỀU thua
   -> kết luận "tối tiểu"      SAI
sự thật:  {x} một mình đã thắng
   -> y và z đều thừa
      \end{dia}
    \end{badb}
  \end{column}
  \begin{column}{0.49\textwidth}
    \lead{muốn chắc thì phải quét \emph{trọn} dàn tập con — $2^{|T|} - 2$ lượt.
          nhưng bảng phủ cắt phần lớn đi:}
    \begin{keyb}[title={chỉ tập con CÒN PHỦ ĐỦ mới cần xét}]
      \begin{dia}
T = {3, 12, 19}
  3 phủ mảnh 1 · 12 phủ mảnh 2 · 19 thừa

tập con thực sự:         7 tập
còn phủ đủ cả hai mảnh:  1 tập -- {3,12}
  {3,19}  mất mảnh 2
  {12,19} mất mảnh 1
  mọi tập một đoạn đều mất một mảnh
      \end{dia}
      \vspace{0.25em}
      và \texttt{c5} tầng 0 \yes{đã đo $\{3,12\}$ rồi} — nên chứng minh tối tiểu
      \hl{tốn thêm 0 lượt}.
    \end{keyb}
    \begin{bx}
      \footnotesize
      \setlength{\tabcolsep}{2pt}
      \begin{tabular}[t]{@{}lrrrrr@{}}
        \toprule
        $|T|$ & 3 & 5 & 6 & 8 & 10\\
        \midrule
        quét hết & \$0,02 & \$0,11 & \$0,24 & \$1,14 & \$5,37\\
        tỉa trước rồi quét & — & \$0,11 & \$0,04 & \$0,07 & \$0,10\\
        \bottomrule
      \end{tabular}
    \end{bx}
  \end{column}
  \end{columns}
\end{frame}

\begin{frame}[fragile]{c7 · ba nhánh kết thúc}
  \framesubtitle{c7}
  \lead{sau \texttt{c6} có đủ số liệu để quyết. ba khả năng, và chỉ khả năng thứ ba
        mới tốn tiền.}

  \begin{columns}[t]
  \begin{column}{0.52\textwidth}
    \begin{bx}[title={ba nhánh}]
      \begin{dia}
(1) đủ, phủ hết mảnh, và tối tiểu
    -> XUẤT T sang tầng b                 xong

(2) chưa đủ  (v(T) < tau_q)
    -> thử ứng viên kế trong danh sách,
       tối đa 3 lần, mỗi lần CHỈ 1 lượt
    -> qua được: về nhánh (1)
       trượt cả 3: sang nhánh (3)

(3) đáp số quá to, hoặc tỉa bị kẹt
    -> CHẠY LẠI đầy đủ cho truy vấn này
      \end{dia}
    \end{bx}
    \begin{keyb}[title={nhánh 2 rẻ hơn tưởng rất nhiều}]
      một lần thử tốn \hl{1 lượt} ($1{,}250$ token); một lần chạy lại đầy đủ tốn
      $113{,}000$ — \yes{90 lần thử mới bằng một lần chạy lại}. ba lần thử vớt
      được ${\sim}66\%$ số ca trượt.
    \end{keyb}
  \end{column}
  \begin{column}{0.44\textwidth}
    \begin{bx}[title={ca «đủ mà không tối tiểu» KHÔNG phải chạy lại}]
      \texttt{c6} đã đo \emph{trọn} dàn tập con dưới $T$, nên tập nhỏ nhất đọc ra
      ngay — dùng lại đúng bảng $x,y,z$ ở slide trước:
      \begin{dia}
v({x,y,z}) = 0,80 >= tau_q   T đủ
   ba cặp con: đều thua
   v({x}) = 0,85 >= tau_q
             -> T KHÔNG tối tiểu
             -> T = {x}
      \end{dia}
      \vspace{0.25em}
      \dd{ca này xảy ra thường xuyên, vì lỗi hay gặp nhất của $g_\theta$ là đề
          xuất tập \no{phình} ra quá cỡ}
    \end{bx}
  \end{column}
  \end{columns}
\end{frame}

\begin{frame}[fragile]{thường có nhiều đáp số đúng — chọn cái nào}
  \framesubtitle{chọn $S^{*}$}
  \lead{ba đoạn cùng nói một ý thì có \emph{ba} tập thắng tối tiểu, mỗi tập chứa
        một đoạn. hai thước đo sau để xếp hạng chúng.}
  \vspace{-0.3em}

  \[
    \beta_i^{+} = \bigl|\{\, S \text{ thiếu } i :
      v(S) < \tau_q \;\wedge\; v(S \cup \{i\}) \ge \tau_q \,\}\bigr|
    \qquad
    \beta_i = \frac{\beta_i^{+} - \beta_i^{-}}{2^{\,n-1}}
    \qquad
    \mathrm{pgi}_i = \bigl|\{\, S \in W : i \in S \,\}\bigr|
  \]
  \dd{$\beta^{+}$: thêm $i$ biến \emph{thua} thành \emph{thắng} ·
      $\beta^{-}$: thêm $i$ làm \emph{mất} chiến thắng ·
      $\mathrm{pgi}$: $i$ có mặt trong bao nhiêu tập thắng tối tiểu}

  \vspace{-0.2em}
  \begin{columns}[t]
  \begin{column}{0.5\textwidth}
    \begin{bx}[title={đếm tay trên đúng bảng $a,b,c$ lúc nãy}]
      \begin{dia}
i = a  {b} 0,18 -> {a,b} 0,94    ĐỔI
       ba ca kia: không           B+ = 1
i = b  {a} 0,12 -> {a,b} 0,94    ĐỔI
       ba ca kia: không           B+ = 1
i = c  không ca nào thua->thắng   B+ = 0
       {a,b} 0,94 thắng
         -> {a,b,c} 0,80 THUA     B- = 1
      \end{dia}
    \end{bx}
  \end{column}
  \begin{column}{0.46\textwidth}
    \begin{keyb}
      \footnotesize
      \begin{tabular}[t]{@{}lrrr@{}}
        \toprule
        & $a$ & $b$ & $c$\\
        \midrule
        $v(\{i\})$ một mình & 0,12 & 0,18 & \no{0,71}\\
        $\beta_i$ & 0,25 & 0,25 & \no{$-0{,}25$}\\
        $\mathrm{pgi}_i$ & 1 & 1 & \no{0}\\
        \bottomrule
      \end{tabular}
    \end{keyb}
    \begin{badb}[title={đây là cả lý do dùng lý thuyết trò chơi}]
      $c$ có \emph{điểm đơn lẻ cao nhất} — xếp theo $v(\{i\})$ sẽ chọn nó. nhưng
      $\beta_c$ \no{âm}: $c$ chưa từng biến thua thành thắng, mà có lần \emph{phá}
      mất một chiến thắng.
    \end{badb}
  \end{column}
  \end{columns}
\end{frame}

\begin{frame}[fragile]{chốt $S^{*}$, và tầng a giao ra cái gì}
  \framesubtitle{chọn $S^{*}$}
  \begin{columns}[t]
  \begin{column}{0.46\textwidth}
    \begin{bx}[title={ba tiêu chí, lọc Pareto rồi gộp}]
      \begin{dia}
càng nhỏ càng tốt    -|S|
cấu trúc             w_b*beta + w_p*pgi
hạng truy hồi        tổng r_i
                     (chỉ dùng phá hoà)
      \end{dia}
      \vspace{0.25em}
      \dd{\textbf{phải chuẩn hoá $\beta$ và pgi riêng rồi mới cộng} —
          $\beta \in [-1,1]$ còn pgi đếm tới hàng chục, cộng thẳng thì pgi
          nuốt mất $\beta$}
    \end{bx}
    \begin{bx}[title={độ tin cậy, miễn phí}]
      nếu chỉ \emph{một} tập thắng rõ thì quy gán chắc. nếu năm tập suýt soát nhau
      thì \no{không chắc} — và ta báo ra con số đó thay vì giấu đi.
    \end{bx}
  \end{column}
  \begin{column}{0.5\textwidth}
    \begin{keyb}[title={đầu ra của tầng a}]
      \begin{dia}
T = S*            <- THỨ DUY NHẤT qua ranh giới

kèm theo, để chẩn đoán về sau:
+ phiên bản g_theta  lô nào gán nhãn
+ đường đi           thẳng · tỉa · thử lại · cứu
+ v(T) - tau_q       vượt ngưỡng bao nhiêu
+ cov(T) / J         phủ mấy trên mấy mảnh
+ W, entropy         mấy đáp số, chắc hay không
+ đoạn beta < 0      ứng viên đoạn gây hại
      \end{dia}
      \vspace{0.25em}
      \hl{nhãn $(q \to T)$ không phụ thuộc kiến trúc model nào} — đào một lần,
      dùng cho mọi lần huấn luyện sau.
    \end{keyb}
  \end{column}
  \end{columns}
\end{frame}

% ════════════════════════════════════════════════════════════════════
\section{tầng b}

\begin{frame}[fragile]{tầng b — câu hỏi hoàn toàn khác}
  \framesubtitle{b · mục tiêu}
  \lead{tầng a nói \emph{đoạn 12 là nguồn}. tầng b hỏi: hệ thống đi \emph{đường nào}
        để tới được đoạn 12? bộ đi tìm ở đây lần theo một đồ thị tri thức, truyền
        tin $L$ bước từ các thực thể trong câu hỏi.}

  \begin{columns}[t]
  \begin{column}{0.46\textwidth}
    \begin{badb}[title={chấm điểm từng cạnh là KHÔNG đủ}]
      điểm từng cạnh cho ra một \emph{đống} cạnh quan trọng, không cho ra một
      \emph{chuỗi}.
      \vspace{0.3em}

      mà lời giải thích đọc được phải có dạng
      \begin{dia}
q -> Forrest Gump -> Zemeckis -> đoạn 12
      \end{dia}
      tức một \hl{dãy cạnh nối tiếp}, mỗi bước ở đúng lượt truyền tin của nó.
    \end{badb}
  \end{column}
  \begin{column}{0.50\textwidth}
    \begin{keyb}[title={ba bước — xương sống của cả phần này}]
      \begin{dia}
b1 ĐO    chấm điểm từng (cạnh, bước)
         -> cạnh nào gánh, và gánh ở bước nào

b2 NỐI   trải phẳng đồ thị theo bước, rồi
         quy hoạch động tìm đường điểm cao nhất
         -> một dãy liền mạch, không phải đống

b3 THỬ   kiểm lại đường đó bằng model thật
         -> nó có thật sự đủ để giải thích không
      \end{dia}
    \end{keyb}
    \begin{bx}[title={đại lượng đem ra chấm}]
      \[ F_B \;=\; \log p(\text{đoạn } D) \]
      \dd{xác suất model gán cho đúng đoạn đích. một lượt chạy cho ra xác suất của
          \emph{mọi} đoạn, nên cả $T$ dùng chung một lượt}
    \end{bx}
  \end{column}
  \end{columns}
\end{frame}
